% cuadro C4_5
\begin{table}[htbp]
\centering
\scriptsize
\caption{Gasto corriente estructural contra su límite máximo, 2026}
\label{cua:C4_5}
\begin{tabular}{p{7.4cm}rr}
\toprule
\textbf{Concepto} & \textbf{2025} & \textbf{2026} \\
\midrule
Gasto corriente estructural (Anexo 2 del decreto) & 3 603 976,7 \\
Limite maximo del gasto corriente estructural & --- \\
Holgura & --- \\
\bottomrule
\end{tabular}
\\[2pt]\begin{minipage}{0.92\textwidth}\scriptsize El límite máximo no es un porcentaje del PIB: se construye del gasto corriente estructural devengado de la Cuenta Pública de 2024 y de los deflactores, con un crecimiento real de 2.05 por ciento contra un PIB potencial de 2.11. El punto de 2025 exige el límite que propuso el CGPE de ese año y no se construye. Fuente: Anexo 2 del proyecto de decreto y CGPE 2026, p. 33. Elaborado por el ITED.\end{minipage}
\end{table}
